% ---------------------------------------------------------------------
%  Conclusion : bilan + ouverture
%  (les limites détaillées et les perspectives vivent au chapitre
%   précédent ; ici, on répond à la question posée en introduction)
% ---------------------------------------------------------------------
\chapter{Conclusion}
\label{chap:conclusion}

L'objectif fixé en introduction n'était pas de découvrir une stratégie
gagnante, mais de construire un protocole capable d'en juger une quelconque.
Ce protocole est désormais construit et éprouvé, et il a rendu son verdict.
Aucune des dix stratégies d'analyse technique ne démontre d'avantage réel une
fois ses conditions d'application vérifiées.

Le résultat central mérite d'être rappelé. Le test de référence déclarait deux
stratégies gagnantes. La seule vérification de la condition d'indépendance a
renversé ces deux verdicts. La première se révélait un faux positif né de la
répétition d'une même information sur des milliers d'opérations, tandis que la
seconde possédait un avantage réel mais trop épisodique dans le temps pour être
retenu. Dans le même temps, le témoin construit pour être gagnant a franchi
chaque étage sans faiblir. Un outil de validation se juge précisément à sa
capacité d'attraper ses propres erreurs tout en reconnaissant ce qui est vrai,
et la démonstration est faite.

Ce verdict négatif n'est pas un échec du travail, il en est le produit
attendu. Rien ne garantissait qu'une des stratégies du banc d'essai possédât
un avantage, et l'analyse technique n'a jamais fait la preuve générale de son
efficacité. Ce que le protocole établit, c'est que l'apparence d'avantage
observée au premier test provenait d'un défaut de méthode, non d'un signal
capté par les stratégies.

Une idée traverse l'ensemble de ce travail. La dépendance des données y a joué
deux rôles opposés. Simultanée, portée par le facteur marché qui explique près
des trois quarts des mouvements, elle était la nuisance qui faussait les
tests, et il a fallu la mesurer puis la neutraliser pour juger proprement.
Retardée, elle aurait au contraire été la seule ressource qu'une stratégie
puisse exploiter, et le marché n'en contient presque pas, conformément à ce
que la théorie financière nomme l'efficience faible. Neutraliser l'une et
chercher l'autre relèvent de la même statistique, vue des deux côtés du même
objet.

Le protocole est ainsi prêt à juger les prochains signaux de l'agent, en
commençant par leur validation temporelle. C'est là, plus que dans le verdict
rendu sur ces onze stratégies, que réside l'apport de ce travail.
